\title{Incremental Value of Daily Workload and Wellness Monitoring Beyond a Personalised Readiness-History Baseline for Next-Calendar-Day Readiness in Elite Women's Football}

\author{Haoyang Cheng\\[0.25em]
\small Wuhan Sport University, Wuhan, China\\
\small ORCID: \href{https://orcid.org/0009-0009-4123-2610}{0009-0009-4123-2610}\\
\small Correspondence: Haoyang Cheng, \href{mailto:1303615018@qq.com}{\texttt{1303615018@qq.com}}}
\date{}

\maketitle

\begin{abstract}
\textbf{Purpose:} Daily workload and wellness monitoring are widely used in elite football, but it is uncertain whether a broader daily panel improves short-horizon readiness forecasts beyond current readiness and the player's own prior forecast errors. We evaluated the incremental forecasting value of a daily workload--wellness panel for next-calendar-day self-reported readiness in elite women's football and examined its stability in a two-team transportability assessment under sequential personalisation.

\textbf{Methods:} This secondary longitudinal forecasting analysis used publicly available SoccerMon subjective-monitoring data from 50 players in two elite women's football teams. The paired primary cohort comprised 14,582 player-day pairs with a complete current-day readiness and wellness profile, observed next-day readiness, and the 7-day load history required by the full model. Predictors available by the end of calendar day $t$ were used to forecast readiness at $t+1$. A current-readiness population model (\Pzero) and a full-monitoring population model (\Pone) were separately tuned in three expanding chronological source-team 2020 folds. During 2021 testing, both models received the same reset, past-only, player-specific residual adjustment. Models were evaluated within teams and bidirectionally across teams. The prespecified primary estimand was \DeltaMAE{} = MAE(\Pone) $-$ MAE(\Pzero), with percentile 95\% player-cluster bootstrap confidence intervals.

\textbf{Results:} The full monitoring model did not reduce MAE in any evaluation setting. \DeltaMAE{} was +0.007 (95\% CI $-$0.005 to +0.023) within Team A, +0.018 ($-$0.005 to +0.047) within Team B, +0.022 ($-$0.005 to +0.053) from Team A to Team B, and +0.004 ($-$0.007 to +0.016) from Team B to Team A; positive values indicate higher error for \Pone. Results were directionally consistent across residual-history windows, a reduced-panel comparison, and a wellness-item-missingness sensitivity analysis conditional on observed current- and next-day readiness.

\textbf{Conclusions:} Among report-complete player-days in these two teams, the evaluated daily workload--wellness panel did not demonstrate stable incremental MAE improvement beyond a personalised readiness-history baseline with identical sequential residual updating. This finding concerns a specific calendar-day prediction task and does not imply that workload or wellness measures are causally irrelevant or lack value for communication, clinical reasoning, or other athlete-management decisions.
\end{abstract}

\textbf{Keywords:} athlete monitoring; women's football; readiness to train; training load; wellness; personalised forecasting; transportability
